% s-036-p036-m001; original PDF 36, 37
\Aphorism{s-036-p036-m001}{22}
Der Mensch als Individuum, als Individualist, fängt heute wieder einmal an langsam heraufzukommen und sich durchsetzen zu wollen; vielleicht glückt’s ihm diesmal besser als früher. Der einzige menschliche Typus, der bisher geduldet und zugelassen wurde, war gerade umgekehrt der typische Mensch, der Mitmensch, das numerirte Mitglied einer Gemeinschaft, das \Greek{ζῶον} \Greek{πολιτικόν}, der vielseitige, gemeinnützige, gemeinverständliche Sittlichkeitsapparat mit Gebrauchsanweisung. Alles gute Gewissen ist zur Zeit in Werthen der Gesellschaft, in socialistischen Werthen angelegt; die Klasse, die Institution, die Heerde, der „Gesammtwille“ fühlt sich unbedingt obenauf — jedes Collectivum weiss sich jedem Individuum gegenüber von vornherein im Recht. In entsprechendem Maasse hat sich das Einzelwesen immer mehr als böses Princip an sich betrachten gelernt; was es auch beginnen mag, es giebt sich zunächst immer Unrecht und lässt sich erst von höheren Instanzen Recht geben: meist von der bestehenden Gesellschaft, oder von einer künftigen oder rein theoretischen Gesellschaftsform, auch im Nothfalle vom „Willen Gottes“, oder dem Interesse der Wissenschaft, oder irgend einer eingebildeten Mission und höheren Bestimmung. Das gesellschaftliche Training hat den Einzelmenschen so übel zugerichtet, sein Wille ist so unsicher und selbst zweiflerisch und innerlich furchtsam geworden, dass er instinctiv nach einer von aussen oder oben ergehenden Rechtfertigung und Bestätigung verlangt, sei es welche sie wolle. So mussten die antiken Religionsstifter jene für uns so peinliche Atmosphäre von Wunderherkunft, übersinnlichen Erkenntnissquellen und göttlicher Inspiration um sich herumlegen — sie mussten es, um Andere und vor allem sich selbst von sich zu überzeugen. So konnte sich der wissenschaftliche Mensch die längste Zeit hin nur als Fortentwicklung des religiösen mit seinem Gewissen vertragen: nur insofern die Erkenntniss der Naturgesetze angeblich zu einer tieferen Gotteserkenntniss führt, nahm man sich die Freiheit, Naturgesetze zu erkennen. So beruft man sich heutzutage auf das allgemeine Wohl, auf den internationalen Gedankenaustausch oder die Förderung grosser synthetischer Ziele, um nicht in der Nacktheit und Beliebigkeit des autonomen Individuums dazustehen —: schon der Abwechslung wegen ist es also höchste Zeit, dass einmal das absolute Ich, die Persönlichkeit, der Einzelmensch mit seinen Ernährungs- und Bildungsgesetzen zu Ehren und vor allem zur Existenz gebracht werde.
\par

